% conclusion.tex -- Data Availability Committee Paper (RU-V6)

\section{Conclusion}

This paper presented Shamir-DAC, a data availability committee design for
enterprise ZK validium systems that achieves information-theoretic data
privacy through Shamir's $(k,n)$-threshold secret sharing. We demonstrated
that the protocol meets all four hypothesis targets: attestation latency
of 175~ms P95 at 500~KB (11$\times$ below the 2-second target), zero
information leakage from individual shares (51/51 privacy tests), reliable
recovery under single-node failure (30/30 tests), and 3.87$\times$ storage
overhead (less than EigenDA's 8$\times$).

The primary contribution is the identification and resolution of a systematic
privacy gap in production validium data availability designs. All deployed
DACs---StarkEx, Polygon CDK, Arbitrum Nova---distribute complete batch data to
every committee member, providing no confidentiality guarantee. For enterprise
deployments handling proprietary operational data, this is a fundamental
limitation. Shamir-DAC demonstrates that information-theoretic privacy
(strictly stronger than computational privacy from encryption) can be
achieved at negligible cost for enterprise batch sizes, requiring only
standard field arithmetic over the BN128 scalar field and ECDSA multi-signature
attestation natively supported by the EVM.

Three directions for future work emerge from this research. First, verifiable
share generation via KZG polynomial commitments or PVSS would remove the
trusted-operator assumption, allowing committee members to independently
verify that their shares are consistent. Second, a network-aware evaluation
with WAN-distributed DAC nodes would quantify the impact of real network
latency on attestation performance. Third, proof-of-custody challenges would
provide cryptographic assurance that committee members retain their shares
beyond the attestation signing event, addressing the data withholding attack
vector.
